\begin{abstract}
强约束低空无人机路径规划需要在障碍物、禁飞区、高度边界和路径平滑性等条件下同时形成可行且稳定的三维轨迹。针对基础进化搜索在窄可行走廊中容易产生不可行个体、低质量绕行和边界贴靠的问题，本文构建一种约束可行性场驱动的自适应进化搜索方法（CVF-AE）。该方法将约束违反信息前移为候选生成阶段的有界方向修正，并通过状态调度、稀疏触发和保守局部接收限制 CVF 的作用范围，避免全种群强制场更新。基于三类合成低空三维场景的实验表明，在所选代表性基线中，CVF-AE 保持 $1.00$ 的可行率，并取得最佳平均排名，在工程轨迹复核标记和额外评价开销之间保持较好折中；消融结果进一步验证了 CVF 候选方向与初始化、稀疏可行性保持机制的互补作用。评价开销分析显示，CVF 分支触发率低于 $1\%$，额外评价次数约为 $1\%$ 量级。结果说明，CVF-AE 可作为静态合成强约束低空路径规划中的均衡型候选生成框架。
\end{abstract}

\noindent\textbf{关键词：} 无人机路径规划；约束优化；进化算法；约束可行性场；启发式优化
